\section*{Chapitre \ref{ch:g1:subtraction} --- Soustraction~: premiers pas}

\begin{solution}{exo:g1:subtraction:1}
$6 - 2 = 4$~; \quad $8 - 5 = 3$~; \quad $10 - 4 = 6$.
\end{solution}

\begin{solution}{exo:g1:subtraction:2}
$9 - 3 = 6$~; \quad $12 - 2 = 10$~; \quad $15 - 4 = 11$.
\end{solution}

\begin{solution}{exo:g1:subtraction:3}
$11 - 8 = 3$ (de $8$ jusqu'\`a $11$)~; \quad
$13 - 9 = 4$~; \quad $16 - 12 = 4$.
\end{solution}

\begin{solution}{exo:g1:subtraction:4}
$2 + 6 = 8$, \quad $6 + 2 = 8$, \quad $8 - 6 = 2$, \quad
$8 - 2 = 6$.
\end{solution}

\begin{solution}{exo:g1:subtraction:5}
$14 - 6 = 8$ (v\'erification~: $8 + 6 = 14$)~; \quad
$17 - 9 = 8$ (v\'erification~: $8 + 9 = 17$).
\end{solution}

\begin{solution}{exo:g1:subtraction:6}
$10 - 3 = 7$~; \quad $10 - 5 = 5$~; \quad $12 - 0 = 12$.
\end{solution}

\begin{solution}{exo:g1:subtraction:7}
$15 - 8 = 7$. Il reste $7$ cerises.
\end{solution}

\begin{solution}{exo:g1:subtraction:8}
\guillemotleft{} $5$ s'\'eloignent \guillemotright{}~: $9 - 5 = 4$
canards. \guillemotleft{} $5$ oies arrivent \guillemotright{}~:
$9 + 5 = 14$ oiseaux. \guillemotleft{} de combien d'ann\'ees plus
\^ag\'e \guillemotright{}~: $9 - 5 = 4$ ans.
\end{solution}

\begin{solution}{exo:g1:subtraction:9}
Famille de $12$, $7$ et $5$~: $12 - 7 = 5$ {\oe}ufs ont \'eclos (et
$5 + 7 = 12$ le v\'erifie).
\end{solution}
